\chapter{基于 SCF-DeepONet 的道岔伤损电磁层析成像模型设计}
\markboth{基于 SCF-DeepONet 的道岔伤损电磁层析成像模型设计}{}

如第二章所述，传统基于灵敏度矩阵的重建方法在铁磁性EMT伤损检测中受限于线性微扰假设，难以处理强非线性与病态性；而标准CNN架构又难以直接处理EMT多激励观测的结构化特征与连续域输出需求。深度算子网络通过在函数空间中学习``结构化观测到连续场分布''的映射，为上述挑战提供了更合适的建模框架。在此基础上，本章面向道岔伤损检测任务，进一步探讨如何从多视角、弱幅值的电磁响应信号中稳定且准确地反演伤损的二维位置与轮廓，是系统实现工程应用的关键算法问题。本章首先从数学与物理层面对道岔伤损检测的 EMT 成像任务进行界定，分析纯数据驱动方法在解决该问题时面临的主要局限；随后，详细介绍本文提出的共享条件特征驱动深度算子网络（Shared-Conditional-Feature Deep Operator Network, SCF-DeepONet）的总体架构、多任务损失函数设计及与之相匹配的多阶段训练策略。

\section{问题描述与深度学习建模思路}

在设计深度学习模型之前，需要首先将道岔伤损的电磁探测过程转化为规范的机器学习问题，并明确模型输入、输出及其对应关系。

\subsection{道岔伤损检测的 EMT 成像任务定义}

本课题面向道岔钢轨及关键部件表面与近表面伤损的无损检测，采用阵列式 EMT 技术进行电磁探测。在实验与仿真论证阶段，本文采用含机加工缺陷的铁磁性试块等效模拟真实道岔伤损。

探测前端由 $4 \times 4$ 阵列排布的 16 个平面线圈构成。检测时系统采用轮流激励方式：对于每一个待测样本，依次选取 16 个线圈中的一个作为激励线圈，其余线圈作为接收线圈，并剔除激励线圈的自接收信号。因此，单个伤损样本在一次完整扫描周期内可获得$16 \times 15 = 240$个复数电压测量值。上述多视角电磁响应数据共同构成深度学习模型的输入特征。

从工程应用目标来看，现场维护人员通常并不需要直接观察复杂的三维电磁场分布，而更关注伤损在检测平面上的位置、范围与大致形状。因此，本文将 EMT 成像任务的主目标定义为：建立从 240 个复数电压响应到二维伤损热力图（2D heatmap）的映射模型。该热力图用于表征伤损在探测平面内的中心位置 $(c_x, c_y)$ 及其几何轮廓分布，从而为后续伤损识别与可视化判读提供依据。

\subsection{数据驱动成像面临的主要挑战}

近年来，端到端卷积神经网络、多层感知机等纯数据驱动方法在图像分割与逆问题求解中取得了较多进展。然而，由于本课题具有明显的电磁物理背景与铁磁材料特性，直接学习``电压响应 $\rightarrow$ 二维热力图''的映射关系仍面临以下挑战：

\begin{enumerate}[label=(\arabic*)]
    \item \textbf{高维多视角输入的特征组织困难。}  
    由 16 次激励产生的 240 个复数电压响应并非相互独立的一维标量集合，而是同时蕴含激励位置、接收位置及其空间相对关系等拓扑信息。若将其简单扁平化后输入常规网络，容易破坏原始测量数据中的空间关联结构，从而限制模型对多视角耦合特征的有效提取。

    \item \textbf{铁磁材料中的磁\cndash 电耦合增强了逆问题复杂性。}  
    道岔钢轨属于典型铁磁性材料。当结构中存在伤损时，不仅局部电导率 $\sigma$ 发生变化，局部磁导率 $\mu$ 也会出现明显扰动。由此导致 EMT 正问题同时表现出磁参数与电参数的空间非均匀性和强非线性耦合特征，进一步加剧了逆问题求解的不适定性。

    \item \textbf{仅依赖二维标签监督易导致``捷径学习''。}  
    若训练阶段仅利用二维矩形伤损标签对网络进行监督，则模型容易学习到电压响应与特定几何轮廓之间的统计捷径，而非底层三维电磁散射机制。此类捷径学习虽然可能在训练分布内取得较高精度，但在面对复杂边界条件、极端工况或未见过的伤损形状（如圆形伤损）时，往往会出现明显的泛化性能下降\upcite{geirhos2020shortcut}。
\end{enumerate}

\subsection{引入辅助场重构任务的基本思想}

针对纯数据驱动方法在物理一致性与泛化能力方面的不足，本文在深度学习建模中引入多任务学习思想：在二维伤损定位主任务之外，增加三维空间矢量磁矢势差场 $\Delta A$ 重构作为辅助任务。

需要明确的是，三维场重构\textbf{并非系统最终部署阶段必须输出的工程指标}，其核心价值主要体现在\textbf{训练阶段}。引入该辅助任务的目的在于：通过要求网络前端的共享特征提取模块同时服务于二维热力图生成与三维差场重构，迫使模型学习电磁响应在空间中的更深层结构信息，从而在一定程度上抑制模型对训练分布中二维几何先验的过拟合，降低其对特定伤损形状模板的依赖，提升模型在极端工况以及分布外样本上的鲁棒性与泛化能力。

从最终部署角度看，系统仅需输出二维伤损热图即可满足工程检测需求。三维差场监督的核心价值不在于直接提升同分布检测精度——实验表明其在矩形样本上的平均 IoU 甚至略低于纯热图模型——而在于以\textbf{训练阶段的物理感知正则化}形式，抑制模型对特定形状模板的过拟合，提升困难样本鲁棒性与分布外泛化能力。更深层次上，该阶段是后续 PINN（物理信息神经网络）强约束嵌入的\textbf{``弱形式预实验''}：通过数据驱动的场--检测耦合，以仿真标签替代复杂 PDE 残差推导，验证物理信息在 EMT 逆问题中的嵌入路径是否可行，并为后期严格麦克斯韦方程微分约束提供经过预热的三维场初始化基础。

\section{数据集构建与样本表示方法}
\label{sec:dataset_construction}

高质量且具有充分物理多样性的数据集，是训练 SCF-DeepONet 模型的前提条件。由于真实道岔现场难以获取大规模、带有精确空间真值的 EMT 成像样本，本文采用有限元仿真方法构建训练与验证数据集。本节将围绕样本生成流程、输入观测表示、输出标签构建以及归一化策略展开说明。

\subsection{有限元仿真模型与样本生成流程}
\label{subsec:comsol_simulation}

本文采用有限元多物理场仿真平台建立 EMT 正问题模型，以铁磁性金属试块等效模拟道岔钢轨或关键部件的基体材料，并在其表面或近表面引入人工伤损，用以模拟实际服役过程中可能出现的疲劳裂纹、剥落等缺陷。样本生成采用多机并行方式，各计算节点通过独立随机种子批量生成矩形缺陷样本，伤损宽度与长度均在 3\,mm 至 18\,mm 之间随机生成，伤损中心坐标根据当前缺陷尺寸动态调整，确保缺陷边缘始终位于探测区域内部。

\begin{figure}[htbp]
    \centering
    \includegraphics[width=0.64\textwidth]{figures/grid_3d.jpg}
    \figcaption{网格剖分示意图}{3D Grid Generation Schematic}
    \label{fig:Grid 3d}
\end{figure}

在网格剖分设置中，为兼顾 EMT 近场敏感区、源场区与远场区域的不同物理特征，本文采用非均匀加密的三维网格结构。具体而言，在EMT传感阵列下方的铁块表面设置了边界层网格以捕捉趋肤效应；在伤损区域及其附近设置了局部加密网格；在远离传感阵列的铁块内部与空气区域则采用较为稀疏的网格以降低计算成本。最终生成完整网格包含数十万个域单元与边界单元，能够在保证仿真精度的前提下实现合理的求解效率。

在每个物理伤损样本求解完成后，仿真程序提取两类核心结果：其一是 16 次激励下所有线圈端口的复数电压响应矩阵 $V_{\text{defect}} \in \mathbb{C}^{16\times16}$；其二是在预先给定的三维固定点云上，对三维磁矢势分量 $A_x$、$A_y$、$A_z$ 进行插值，得到形状为 $16\times N\times3$ 的复数场响应数据，其中 $N$ 表示固定评估点数。最终，单个物理样本被保存为一个包含伤损标签、端口电压响应以及三维场响应的独立数据文件。

\begin{figure}[ht]
    \centering
    \subfigure[\shortstack{(a) X-Z 侧视图 \\ {\zihao{5}X-Z Side View}}]{\includegraphics[width=0.46\textwidth]{figures/Fig2_XZ_Side_View.jpg}}
    \hspace{2mm}
    \subfigure[\shortstack{(b)X-Y 俯视图 \\ {\zihao{5}X-Y Top View}}]{\includegraphics[width=0.46\textwidth]{figures/Fig3_XY_Top_View.jpg}} \\[-6pt]
    \subfigure[\shortstack{(c)三维全局视角 \\ {\zihao{5}Overall 3D Global View}}]{\includegraphics[width=0.62\textwidth]{figures/Fig1_3D_Global_View.jpg}}
    \figcaption{采样点云图分布}{Sampling point cloud distribution}
    \label{fig:Evaluation Points Distribution}
\end{figure}

为了兼顾 EMT 近场敏感区、源场区与远场区域的不同物理特征，本文并未采用均匀体素网格作为三维场采样点，而是预先生成一组全局固定的非均匀三维评估点云，其分布如图 \ref{fig:Evaluation Points Distribution}所示。该点云总数为 50000 个，按空间物理意义划分为四个区域：其一为核心敏感区，共 40000 个点，覆盖铁块内部及伤损附近区域；其二为源场区，共 5000 个点，用于刻画线圈本体及其上方空气间隙中的激励源分布；其三为远场空气区，共 2000 个点，用于表征空气中的衰减场趋势；其四为深层/外围铁块区，共 3000 个点，用于提供深层金属内部的边界与弱场信息。四部分点云最终拼接后随机打乱顺序，形成全局固定评估点集。

\subsection{输入响应数据构建}
\label{subsec:input_construction}

EMT 系统前端由 $4\times4$ 阵列排布的 16 个平面线圈构成。对任意一个伤损样本，系统采用轮流激励模式依次驱动 16 个线圈中的一个，其余线圈作为接收端采集感应电压。由于自激励自接收信号中包含极强的一次背景场分量，且其量级远大于伤损所引起的微小扰动，因此在构造网络输入时，本文剔除了所有自激励自接收项，仅保留每次激励下其余 15 个接收通道的响应。由此，单个物理样本在一次完整扫描中共获得240个复数电压测量值。

然而，本文并未将这 240 个复数值简单地扁平化为一维向量，而是将其组织为一个具有明确几何语义的张量表示。具体而言，首先以无伤损基准模型输出的端口电压 $V_b$ 作为参考，将每个伤损样本对应的电压矩阵 $V_d$ 转化为相对差分响应

\begin{equation}
\label{eq:V_normlization}
V_{\delta,\mathrm{norm}}=\frac{V_b-V_d}{|V_b|+\varepsilon}
\end{equation}

其中 $\varepsilon$ 为防止分母过小而引入的微小常数。该相对差分形式能够在保留伤损扰动信息的同时，显著削弱大背景场幅值差异对网络训练的影响。

在此基础上，本文进一步为每一组``激励线圈—接收线圈对''构造了几何拓扑特征。设激励线圈中心坐标为 $(x_t,y_t)$，接收线圈中心坐标为 $(x_r,y_r)$，则附加几何特征包括激励坐标、接收坐标、相对位移 $(\Delta x,\Delta y)$ 以及二者欧氏距离 $d$，共 7 个量。所有几何量均采用与阵列尺度相匹配的固定系数进行归一化处理。最终，每个激励视角下的 15 个接收响应被表示为一个 $15\times9$ 的特征矩阵，其中前 2 维为复数差分电压的实部与虚部，后 7 维为对应的几何拓扑特征；16 个激励视角共同构成形状为$16\times15\times9$的输入张量 $v_{\mathrm{pair}}$。这种表示方式保留了多激励、多接收以及空间几何关系之间的结构信息，比简单扁平化输入更适合后续共享观测编码器的学习。

\subsection{输出标签构建}
\label{subsec:output_construction}

SCF-DeepONet 采用二维检测与三维场重构相结合的多任务学习框架，因此每个样本需要同时构建二维热图标签和三维差场标签。

\textbf{（1）二维伤损热图标签：}  
二维热图标签用于监督模型的伤损检测任务。对于每个矩形伤损样本，使用其真实几何参数$(c_x,c_y,w,l)$在二维探测平面上定义伤损区域，并据此生成二值掩膜：伤损内部取值为 1，背景区域取值为 0。需要说明的是，在实际训练过程中，本文并不是在整张规则网格上直接进行像素级监督，而是采用分层采样方式动态生成二维训练点集。具体地，在每个样本的训练过程中，从伤损内部、伤损边界附近以及全局背景区域分别采样二维点，并根据点是否落入真实伤损区域来赋予监督标签。该策略能够在保证训练效率的同时，使网络更加关注伤损边界与局部几何细节。此外，在后续的零样本泛化实验中，只需改变掩膜的几何生成方程，即可将矩形伤损标签替换为其他形状的伤损标签。

\textbf{（2）三维差场重构标签：}  
三维场标签用于引导模型学习伤损对空间电磁散射结构的影响。本文从有限元软件中提取固定三维点云上的磁矢势分量$A=(A_x,A_y,A_z)$作为基础场量，并定义差分磁矢势场
\begin{equation}
\label{eq:delta_A}
\Delta A = A_{\mathrm{defect}}-A_{\mathrm{baseline}}
\end{equation}
作为场重构任务的监督目标。为便于神经网络处理复数矢量场，本文将每个空间点上的复数三分量磁矢势拆分为六个实数通道，即
\[
(\Re(A_x),\Im(A_x),\Re(A_y),\Im(A_y),\Re(A_z),\Im(A_z))
\]
因此，对于任意一个物理样本，16 次激励下的三维差场标签最终可表示为形状为$16\times50000\times6$的张量。除差场标签外，数据集中还同时保存了无伤损状态下的基准场$A_0$，并采用相同方式拆分与归一化。这样做的目的在于：一方面便于 Stage 2 的差场重构训练；另一方面，也为 Stage 4 中构造总场
\begin{equation}
\label{eq:A_total}
A_{\mathrm{total}}=A_0+\Delta A
\end{equation}
并进一步引入偏微分方程弱约束提供必要的数据支撑。

\subsection{数据归一化与坐标表示}
\label{subsec:data_normalization}

由于输入电压、三维场量以及空间坐标在量纲和数值范围上存在显著差异，若直接将其用于网络训练，容易导致优化过程不稳定。因此，本文针对不同类型数据采用了分层归一化策略。

\textbf{（1）二维表面坐标与三维场坐标的分离归一化：}  
考虑到二维热图任务与三维场任务对空间范围的关注尺度并不相同，本文分别构建了二维表面坐标归一化与三维场坐标归一化两套体系。对于二维热图相关的表面投影坐标，采用敏感区半宽
\[
\mathrm{MAX\_SURF\_XY}=32.0\ \mathrm{mm}
\]
进行归一化；对于三维场重构相关坐标，则采用
\[
\mathrm{MAX\_FIELD\_XYZ}=[130,\,90,\,90]\ \mathrm{mm}
\]
分别对 $x$、$y$、$z$ 三个方向进行缩放。上述两套坐标系在数据集打包阶段被同时保存，并在网络的二维热图查询与三维场查询中分别使用。

\textbf{（2）电压响应归一化：}  
对于输入电压，本文采用相对差分归一化，即以前述式\ref{eq:V_normlization}作为网络输入。

\textbf{（3）三维差场与基准场归一化：}  
对于三维磁矢势场，本文通过全局统计量构造统一尺度因子。具体做法为：首先对全部伤损样本的差场幅值 $|\Delta A|$ 分别计算其 99\% 分位数，再取这些分位数的中位数作为全局缩放系数 $A_{\mathrm{scale}}$。随后，差场 $\Delta A$ 与基准场 $A_0$ 均使用同一 $A_{\mathrm{scale}}$ 进行归一化。该策略避免了少量异常极值对整体尺度造成过度影响，并使不同样本、不同激励下的场量均被映射到相近数量级。

上述输入输出与监督信号的定义与构建方法如表\ref{tab:input_output_definitions}所示。

\begin{table}[ht]
    \centering
    \tabcaption{SCF-DeepONet 输入、输出与监督信号定义}{Definitions of Inputs, Outputs and Supervision Signals in SCF-DeepONet}
    \vspace{6bp}
    \begin{threeparttable}
    \begin{tabularx}{0.95\textwidth}{
        >{\centering\arraybackslash}p{0.18\textwidth}
        >{\centering\arraybackslash}p{0.18\textwidth}
        >{\centering\arraybackslash}p{0.18\textwidth}
        >{\raggedright\arraybackslash}X
    }
        \toprule
        类别 & 符号 & 维度/形式 & 说明 \\
        \midrule
        输入观测 & $v_{\delta,\mathrm{norm}}$ & 结构化响应张量 & 由多激励条件下的接收线圈差分电压构成的多视角观测输入。 \\
        二维查询坐标 & $(x,y)$ & 连续坐标 & 二维热图分支的坐标查询输入。 \\
        三维查询坐标 & $(x,y,z)$ & 连续坐标 & 三维场分支的坐标查询输入。 \\
        激励索引 & $e$ & 离散索引 & 指定当前三维场查询对应的激励线圈编号。 \\
        二维监督标签 & heat\_target & 二值掩膜/点标签 & 由真实伤损几何参数生成的二维热图监督。 \\
        三维监督标签 & $A_{\delta,\mathrm{norm}}$ & 连续场值 & 归一化后的三维磁矢势差场监督。 \\
        基准场 & $A_{0,\mathrm{norm}}$ & 连续场值 & 用于 Stage 4 中构造总场的无伤损基准场。 \\
        PDE 物理约束 & $L_{\mathrm{PDE,rel}}$ & 相对残差 & 基于总场、材料分布和控制方程构造的弱物理约束项。 \\
        \bottomrule
    \end{tabularx}
    \end{threeparttable}
    \label{tab:input_output_definitions}
\end{table}

\section{SCF-DeepONet 网络总体架构}
\label{sec:network_architecture}

针对道岔伤损 EMT 成像中多激励输入维度高、空间耦合关系复杂以及检测与场重构并行建模的需求，本文在经典 DeepONet 思想的启发下，构建了一种面向铁磁性伤损检测的扩展多任务网络——共享条件特征驱动深度算子网络。与标准 DeepONet 侧重``单一算子映射''的 Branch--Trunk 结构不同，本文所提出的网络在共享观测编码基础上，同时耦合二维伤损热图生成与三维差场重构两个任务分支，从而形成面向 EMT 检测场景的多任务算子学习框架。

SCF-DeepONet 网络总体架构如图 \ref{fig:network_architecture} 所示。从程序实现上看，最终网络主控类由五个核心子模块组成：配对编码器（pair\_encoder）、激励池化模块（pool）、特征精炼器（refiner）、二维热图分支（heatmap\_trunk）以及三维场分支（field\_trunk）。其中，共享观测编码模块负责将输入观测张量编码为可供二维与三维分支共享使用的全局与局部条件特征；二维分支完成伤损热图的连续坐标查询；三维分支则在二维伤损先验引导下重构连续空间中的磁矢势差场 $\Delta A$。

\begin{figure}[ht]
    \centering
    \includegraphics[width=0.96\textwidth]{figures/SCF_DeepONet_Structure.png}
    \figcaption{SCF-DeepONet 总体网络架构图}{Overall Architecture of SCF-DeepONet}
    \label{fig:network_architecture}
\end{figure}

\subsection{总体框架设计}
\label{subsec:overall_framework}

经典深度算子网络旨在学习从一个输入函数到输出函数的映射算子。设输入函数为 $u \in \mathcal{U}$，输出函数所在空间为 $\mathcal{Y}$，则待学习算子 $G: \mathcal{U} \to \mathcal{Y}$ 在查询点 $y \in \Omega$ 处的值可通过 Branch--Trunk 结构逼近：
\begin{equation}
\label{eqn:deeponet_general}
\hat{G}(u)(y)
=
\sum_{k=1}^{p} b_k(u) \cdot t_k(y) ,
\end{equation}
其中 $b_k(u) = \mathrm{Branch}_k(u)$ 为 Branch 网络对离散观测 $u$ 编码得到的 $p$ 维隐特征，$t_k(y) = \mathrm{Trunk}_k(y)$ 为 Trunk 网络在查询坐标 $y$ 处求值得到的 $p$ 维基函数响应。该结构将观测编码与空间求值解耦，使输出可在任意连续坐标上查询，摆脱了固定输出网格的限制。

本文提出的 SCF-DeepONet 在经典 DeepONet 基础上进行了面向 EMT 检测任务的扩展。设输入观测为结构化电压响应张量 $v_{\mathrm{pair}}$，二维伤损热图分支和三维差场分支可分别写为
\begin{align}
\label{eqn:scfe_deeponet_heat}
\hat{G}_{\mathrm{heat}}(v)(x,y)
&= \mathrm{HeatmapTrunk}\bigl(
    f_{xy},\, z_g,\, \gamma(x,y)
\bigr) , \\
\label{eqn:scfe_deeponet_field}
\hat{G}_{\mathrm{field}}(v)(x,y,z)
&= \mathrm{FieldTrunk}\bigl(
    f_{xyz},\, h_{\mathrm{feat}},\, h_{\mathrm{prob}},\, z_g,\, z_e,\, \gamma(x,y,z)
\bigr) ,
\end{align}
其中 $\gamma(\cdot)$ 为 Fourier 特征映射\citep{tancik2020fourier}，$f_{xy}$ 和 $f_{xyz}$ 分别为二维和三维坐标对应的局部特征，$h_{\mathrm{feat}}$ 和 $h_{\mathrm{prob}}$ 为二维热图分支提供的伤损先验特征与概率，$z_g$ 和 $z_e$ 分别为全局条件特征与激励条件特征。

网络模型的输入为形状为 $16\times15\times9$ 的结构化观测张量 $v_{\mathrm{pair}}$。网络首先通过共享观测编码模块提取两个层次的条件特征：其一为可在二维平面上双线性采样的局部特征图 $F_{map}$，其二为表征全局观测状态的特征向量 $z_g$。此外，编码模块还保留每个激励对应的独立特征向量 $z_e$，用于刻画当前激励视角下的专属信息。

在此基础上，二维热图分支接收任意二维查询坐标 $(x,y)$，结合由 $F_{map}$ 采样得到的局部特征 $f_{xy}$ 以及全局特征 $z_g$，输出该位置对应的伤损概率；三维场分支则接收任意三维查询坐标 $(x,y,z)$、当前激励索引以及由二维分支提供的伤损先验特征，在连续空间中输出对应位置的三维复数差场分量 $\Delta A$。与简单的``双头网络''不同，本文三维场分支并不是完全独立于二维分支运行，而是显式调用二维热图分支，将其产生的深层特征与伤损概率作为三维场预测的重要条件输入。正是这种从``共享编码 $\rightarrow$ 二维检测 $\rightarrow$ 三维场重构''的层级耦合关系，使得三维辅助任务能够反向约束二维检测相关特征的物理一致性。 

\subsection{共享观测编码模块}
\label{subsec:shared_encoder}

共享观测编码模块承担了从多激励、多接收电压响应中提取判别性特征的任务，其整体流程可表示为
\[
v_{\mathrm{pair}}
\longrightarrow
\mathrm{PairEncoder}
\longrightarrow
\mathrm{Pool}
\longrightarrow
\mathrm{Refiner}
\longrightarrow
\{F_{map},z_g,z_e\}.
\]

首先，配对编码器（Pair Encoder）以``激励线圈--接收线圈对''为基本处理单元，对每一个线圈对的 9 维输入特征进行非线性映射，获得局部观测特征。由于输入观测本身以 $16\times15$ 的配对形式组织，因此该模块能够在不破坏激励--接收结构的前提下完成第一层特征提取。随后，激励池化模块（Pool）对每个激励视角下的 15 个接收响应进行聚合，形成每个激励对应的紧凑表示 $z_e$；这些激励特征在 batch 维度上重新组织后，再送入特征精炼器（Refiner），进一步生成二维局部特征图 $F_{map}$ 和全局特征向量 $z_g$。共享编码器最终同时输出全局特征 $z_g$、激励特征集合以及二维局部特征图
其中，$F_{map}$ 的作用可以理解为``由多激励观测反推出的二维潜在伤损特征场''，后续二维热图分支和三维场分支都需要在给定空间坐标处从该特征图中进行采样；$z_g$ 则用于提供样本级全局条件；而 $z_e$ 则在三维场预测中用于区分不同激励视角下的空间场差异。共享观测编码模块的输出形式为后续 2D/3D 双分支协同学习奠定了统一表征基础。

\subsection{二维伤损热图分支}
\label{subsec:heatmap_trunk}

二维伤损热图分支（Heatmap Trunk）承担了本文最核心的工程任务，即在检测平面上估计伤损位置与形状分布。与传统卷积解码器直接输出整幅图像不同，本文采用基于连续坐标查询的隐式表示方式：对于任意给定的二维归一化坐标 $(x,y)$，首先在共享特征图 $F_{map}$ 上通过双线性插值得到局部特征 $f_{xy}$，再将坐标本身、局部特征 $f_{xy}$ 以及全局条件 $z_g$ 共同输入热图分支，输出该位置的伤损响应。

从结构上看，热图分支首先对二维坐标进行 Fourier 特征映射，再与局部特征拼接，经若干层 FiLM 调制的全连接网络后\citep{perez2018film}，分别输出三类中间量：其一为深层隐特征 heat\_feat，其二为最终逻辑值heat\_logits，其三为经 Sigmoid 激活后的伤损概率heat\_prob。其中，heat\_feat 不仅是热图分支的中间表征，也是连接二维检测与三维场预测的关键桥梁；heat\_logits 和heat\_prob 则分别用于损失计算与伤损概率解释。

这种设计具有两方面优势：其一，采用坐标查询方式后，二维热图分支不再依赖固定分辨率的卷积解码器，而可以在任意给定二维点集上进行预测，便于训练时采用分层采样点监督以及验证时在规则网格上高分辨率评估；其二，热图分支中保留下来的深层特征与概率响应，能够作为``伤损位置先验''继续传递给三维场分支，从而建立从二维检测到三维场重构的特征级耦合关系。

\subsection{三维差场重构分支}
\label{subsec:field_trunk}

三维差场重构分支（Field Trunk）是本文引入的关键辅助任务模块，其目标是预测连续三维空间中任意位置的磁矢势差场
\[
\Delta A=(\Delta A_x,\Delta A_y,\Delta A_z).
\]
考虑到电磁散射场不仅取决于空间坐标本身，还与当前激励线圈位置及伤损分布密切相关，三维场分支采用了多条件联合输入的设计。

对于任意一个三维归一化查询点 $(x,y,z)$，网络首先根据当前样本对应的激励索引，从共享编码结果中选取当前激励的特征向量 $z_e$；然后将三维点在物理坐标意义上投影到检测平面，并在二维特征图 $F_{map}$ 上采样得到该点正上方对应的局部二维特征；在此基础上，网络进一步调用二维热图分支，在投影点位置处获得该点对应的heat\_feat 和heat\_prob。最后，将三维坐标的 Fourier 特征、局部特征 $f_{xy}$、二维先验特征heat\_feat、二维伤损概率 heat\_prob、全局条件 $z_g$ 以及激励条件 $z_e$ 一并送入场分支，输出 6 通道实数形式的差场预测结果，对应于
\[
(\Re(\Delta A_x),\Im(\Delta A_x),\Re(\Delta A_y),\Im(\Delta A_y),\Re(\Delta A_z),\Im(\Delta A_z)).
\]
这种多条件联合输入方式使得三维场分支能够充分利用二维检测先验、空间几何信息以及全局观测状态，从而在连续三维空间中实现更准确的差场预测。

特别需要指出的是，三维场分支并不是简单地与二维分支``并列存在''，而是在前向传播过程中显式调用二维热图分支，以获得伤损位置相关的中间特征与概率信息，再在此基础上预测三维差场。也就是说，二维分支在本网络中不仅承担最终检测输出，还充当了三维场分支的重要条件先验。这种``2D 先验引导 3D 场''的层级设计，是本文区别于普通多头网络的关键所在：它使得三维场重构任务能够在训练阶段迫使共享编码器保留与真实三维散射机制一致的底层信息，从而为后续引入基于总场的偏微分方程弱约束微调奠定了网络结构基础。

\section{多任务损失函数设计}
\label{sec:loss_function}

SCF-DeepONet 同时包含二维伤损热图预测与三维差场重构两个任务分支。由于二者在监督形式、数据分布以及优化难度上存在明显差异，若采用单一固定损失对整个训练过程进行端到端优化，容易导致主任务与辅助任务之间的梯度竞争，从而影响模型稳定收敛。因此，本文首先分别定义二维检测损失与三维场重构损失，再结合后续的分阶段训练策略，构建阶段相关的优化目标。

\subsection{热图检测损失}
\label{subsec:heatmap_loss}

二维伤损热图分支的本质是一个连续坐标点上的二分类任务，即判断给定二维坐标 $(x,y)$ 是否位于伤损区域内部。为兼顾点级分类稳定性与前景区域重叠度优化，本文将热图损失定义为二元交叉熵损失与 Dice 损失的加权组合\cite{milletari2016v}。

首先，设第 $i$ 个二维采样点的真实标签为 $y_i\in\{0,1\}$，网络输出的原始逻辑值为 $\hat{l}_i$，则基于 Sigmoid 激活的二元交叉熵损失可写为

\begin{equation}
\label{eq:L_BCE}
L_{BCE}
=
-\frac{1}{N_{2D}}
\sum_{i=1}^{N_{2D}}
\left[
y_i\log \sigma(\hat{l}_i)
+
(1-y_i)\log\left(1-\sigma(\hat{l}_i)\right)
\right]
\end{equation}

其中，$N_{2D}$ 表示单次训练迭代中参与监督的二维采样点总数，$\sigma(\cdot)$ 为 Sigmoid 函数。实际实现中，本文采用数值稳定性更好的 \verb|BCEWithLogitsLoss| 直接计算该项损失。

然而，在 EMT 伤损检测任务中，伤损区域在整个二维探测平面内所占比例通常较小，存在明显的前景--背景不平衡问题。若仅使用 BCE 损失，模型容易偏向于将大多数采样点判定为背景区域。为缓解这一问题，本文进一步引入 Dice 损失，以直接约束预测区域与真实区域之间的重叠程度。设 $\hat{p}_i=\sigma(\hat{l}_i)$ 为第 $i$ 个采样点属于伤损区域的预测概率，则 Dice 损失定义为

\begin{equation}
\label{eq:L_Dice}
L_{Dice}
=
1-
\frac{
2\sum_{i=1}^{N_{2D}} \hat{p}_i y_i + \varepsilon
}{
\sum_{i=1}^{N_{2D}} \hat{p}_i
+
\sum_{i=1}^{N_{2D}} y_i
+
\varepsilon
}
\end{equation}

其中，$\varepsilon$ 为防止分母为零而引入的极小常数，本文取 $\varepsilon=10^{-6}$。最终，二维热图分支的总损失定义为

\begin{equation}
\label{eq:L_heat}
L_{heat}=L_{BCE}+\lambda_{dice}L_{Dice}
\end{equation}

在本文实验中，Dice 损失权重系数取为 $\lambda_{dice}=0.5$。

\subsection{三维场重构损失}
\label{subsec:field_loss}

三维差场重构分支的目标是在连续三维空间中预测磁矢势差场$\Delta A$。考虑到实际训练中复数三分量场已被拆分为六通道实数表示，因此三维场分支可视为一个多输出连续回归问题。本文采用均方误差（Mean Squared Error, MSE）作为三维场重构损失，其定义为

\begin{equation}
\label{eq:L_A}
L_A
=
\frac{1}{N_{3D}}
\sum_{j=1}^{N_{3D}}
\left\|
\Delta \hat{A}_j-\Delta A_j
\right\|_2^2
\end{equation}

其中，$N_{3D}$ 表示单次迭代中参与监督的三维采样点总数，$\Delta \hat{A}_j$ 与 $\Delta A_j$ 分别为第 $j$ 个三维采样点处的预测差场与真实差场。

需要强调的是，$L_A$ 并不是部署阶段的最终工程指标，而是训练阶段的辅助监督信号。其主要作用在于促使共享编码器在学习二维伤损检测的同时，保留与三维电磁散射结构相关的判别性表征，从而增强模型的物理一致性与泛化能力。

\subsection{分阶段联合优化目标}
\label{subsec:joint_loss}

由于二维热图检测任务与三维差场重构任务在监督形式、收敛速度及梯度尺度上存在明显差异，本文并未采用单一固定权重在整个训练过程中同时优化全部目标，而是结合模型训练进程，构建阶段相关的优化目标。

在 Stage 1 中，模型首先以二维伤损检测为唯一目标进行预训练，其优化目标定义为

\begin{equation}
\label{eq:L_stage1}
L^{(1)} = L_{heat}
\end{equation}

当热图分支已具备初步定位能力后，在 Stage 2 中引入三维差场监督，构成联合训练目标：
\begin{equation}
\label{eq:L_stage2}
L^{(2)} = \frac{1}{16}L_{heat} + w_A L_A
\end{equation}
其中，系数 $1/16$ 用于平衡二维热图损失与三维场损失在数值尺度上的差异，$w_A$ 为三维场监督项的权重系数。实际训练中，为减弱三维辅助任务在联合训练初期对共享表征的冲击，$w_A$ 采用逐步增大的 warmup 方式进行调节\citep{goyal2017accurate}。

在 Stage 3 中，为进一步提升二维检测精度，优化目标重新退化为仅包含热图损失的形式：
\begin{equation}
\label{eq:L_stage3}
L^{(3)} = L_{heat}
\end{equation}
该阶段的作用是：在尽量不破坏前期联合训练所形成共享表征的前提下，对热图输出层进行小幅校准。

需要说明的是，后续 Stage 4 所涉及的 PDE 弱物理约束已超出纯监督学习框架，其优化目标将与偏微分方程相对残差项共同构成。考虑到该部分不仅涉及损失函数变化，还涉及总场构造、材料参数场近似以及数值稳定化处理，因此本文将其放在第四章中结合物理约束引入过程统一讨论。

由此可见，SCF-DeepONet 的训练并非依赖固定加权的静态多任务优化，而是通过``先检测、再联合、后精修''的阶段化目标切换，使模型能够逐步完成二维定位能力建立、三维物理表征学习以及检测边界校准等不同子目标。

\section{四阶段训练策略}
\label{sec:curriculum_training}

SCF-DeepONet 同时包含二维伤损检测、三维差场重构以及后续的物理一致性增强等多个子目标。由于这些目标在优化难度、收敛阶段以及对共享特征空间的作用方式上并不一致，若直接采用统一的端到端训练方式，容易导致主任务与辅助任务在训练早期相互牵制，使模型过早收敛到次优解。为此，本文设计了一种如图\ref{fig:four_stage_training}所示的由易到难、逐层递进的四阶段训练策略，通过分阶段切换训练目标并配合不同的参数冻结方式，使模型依次获得二维定位能力、三维场建模能力、热图边界校准能力以及物理一致性增强能力\citep{wang2021survey}。

\begin{figure}[htbp]
    \centering
    \includegraphics[width=0.9\textwidth]{figures/SCF_DeepONet_four_stage_curriculum.png}
    \figcaption{SCF-DeepONet 四阶段训练流程图}{Four-Stage Training Process of SCF-DeepONet}
    \label{fig:four_stage_training}
\end{figure}

\subsection{Stage 1：热图预训练}
\label{subsec:stage1}

训练初期的首要目标，是使模型建立从多激励观测响应到二维伤损分布的基础映射关系。因此，在 Stage 1 中，网络被视为一个纯二维检测模型进行训练。此阶段仅利用二维热图标签对共享观测编码模块和热图分支进行监督，使模型优先学习与伤损位置和轮廓最直接相关的观测特征。采用这一训练顺序的主要原因在于：若在共享特征尚未形成基本二维判别能力之前即引入三维连续场监督，往往会使模型在训练初期承受过大的优化负担，从而影响主任务的稳定收敛。通过先完成热图预训练，网络能够快速建立基础的空间定位能力，并为后续多任务联合训练提供较优的参数初始化状态。

\subsection{Stage 2：热图--场联合训练}
\label{subsec:stage2}

在模型已具备基础二维定位能力后，训练进入 Stage 2，即二维热图与三维差场的联合训练阶段。此时，三维场分支被正式纳入优化过程，共享观测编码模块需要同时服务于二维伤损检测与三维差场预测两个任务。

这一阶段是整套训练流程中最关键的部分。与 Stage 1 相比，网络不再只需学习``观测响应——二维轮廓''之间的直接对应关系，而必须同时捕捉与三维电磁散射结构相关的特征信息。由于二维与三维任务共享前端编码模块，因此三维差场监督会反向约束共享特征空间，使其从单纯面向矩形热图拟合的表示，逐步演化为兼具几何判别性和物理一致性的共享表征。

从方法设计角度看，Stage 2 的意义不只在于提升三维场分支本身的重构精度，更在于通过三维辅助任务改善前端共享表征的质量。后续在困难样本及分布外样本上的泛化性能提升，本质上正来源于这一阶段多任务联合训练对共享特征空间的重塑作用。

\subsection{Stage 3：热图精修}
\label{subsec:stage3}

尽管 Stage 2 的联合训练增强了模型的三维物理表达能力，但多任务训练不可避免地会带来一定程度的梯度竞争，使二维热图分支的最终决策边界出现轻微模糊或偏移。为进一步逼近二维检测任务的最优性能，本文在联合训练之后引入 Stage 3 热图精修阶段。

该阶段从 Stage 2 中检测性能最优的模型出发，采用``主干冻结、头部微调''的策略：冻结共享观测编码模块、三维场分支以及热图分支的大部分参数，仅保留热图输出头为可训练状态，并以较小学习率继续进行优化。这样做的目的是，在不显著破坏 Stage 2 已形成共享物理表征的前提下，对最终热图决策边界进行专门校准，从而进一步提升二维检测精度。

因此，Stage 3 的本质并不是重新训练整个模型，而是对已经形成的共享表示进行低幅度检测校准。这一阶段对应于本文监督学习主线中的最终检测精修步骤，也是后续构建最优多任务检测基线的重要来源。

\subsection{Stage 4：PDE 弱物理微调}
\label{subsec:stage4}

前三个阶段本质上仍属于监督式数据驱动训练范畴。尽管经过 Stage 2 和 Stage 3 后，模型已经分别获得了较好的三维场表达能力与二维检测性能，但三维场预测仍可能出现局部不符合电磁学规律的非物理畸变。为进一步改善三维场分支的物理一致性，本文在最后引入 Stage 4 的 PDE 弱物理微调阶段。

需要强调的是，该阶段的核心作用对象是\textbf{三维场分支}，其主要目标是在不破坏二维检测性能的前提下，对场重构结果施加基于麦克斯韦方程的弱正则约束。由于 PDE 弱约束仅微调三维场分支参数，二维热图分支和共享编码模块均处于冻结状态，因此\textbf{该阶段并不直接改变或提升二维伤损检测指标}。它的价值在于验证：在已经具备较好检测性能的基础上，能否进一步使网络内部的三维场预测满足底层电磁规律，从而为模型的可解释性与物理可靠性提供额外保障；同时，作为 Stage 2 预实验的后续强化，为物理信息神经网络在 EMT 领域的深入应用提供完整经验。

与 Stage 2 的三维场辅助监督相比，Stage 4 的 PDE 弱约束不再是数据驱动的近似，而是以控制方程微分残差形式施加的严格物理正则；但二者的递进关系不变：前者完成物理嵌入的路径验证与模型预热，后者在此基础上完成局部精修。

综上，四阶段训练策略设置总结为表\ref{tab:curriculum_training}。其核心并不在于简单增加训练步骤，而在于依据不同阶段的主要任务，将``检测能力建立''``共享表征重塑''``热图边界校准''与``场分支物理增强''有机拆分并按序实施。该策略使 SCF-DeepONet 能够在保证二维伤损检测主任务稳定收敛的同时，逐步吸收三维场辅助监督与 PDE 物理约束带来的收益，从而形成兼顾检测性能、场重构质量与物理一致性的整体模型。

\begin{table}[ht]
    \centering
    \tabcaption{SCF-DeepONet 四阶段课表训练策略设置}{Curriculum Training Settings for SCF-DeepONet}
    \vspace{6bp}
    \begin{threeparttable}
    \begin{tabularx}{0.95\textwidth}{
        >{\centering\arraybackslash}p{0.16\textwidth}
        >{\centering\arraybackslash}p{0.18\textwidth}
        >{\centering\arraybackslash}p{0.18\textwidth}
        >{\raggedright\arraybackslash}X
    }
        \toprule
        阶段 & 训练目标 & 优化目标 & 参数更新策略 \\
        \midrule
        Stage 1 & 热图预训练 & $L^{(1)}=L_{heat}$ & 训练共享观测编码模块与二维热图分支，建立基础二维定位能力。 \\
        Stage 2 & 热图--场联合训练 & $L^{(2)}=\frac{1}{16}L_{heat}+w_A L_A$ & 联合训练共享编码模块、二维热图分支和三维场分支，引入三维差场辅助监督。 \\
        Stage 3 & 热图精修 & $L^{(3)}=L_{heat}$ & 冻结共享编码模块与三维场分支，仅微调热图输出头，校准二维检测边界。 \\
        Stage 4 & PDE 弱物理微调 & $L^{(4)} \approx L_A+\lambda_{pde}L_{\mathrm{PDE,rel}}$ & 回滚至 Stage 2 场最优模型，冻结共享编码模块与二维热图分支，仅微调三维场分支。 \\
        \bottomrule
    \end{tabularx}
    \end{threeparttable}
    \label{tab:curriculum_training}
\end{table}

\section{本章小结}

本章围绕道岔伤损 EMT 检测任务，完成了 SCF-DeepONet 模型的构建。首先，结合高导电率、高磁导率铁磁性目标的检测特点，明确了以多视角复数差分电压响应为输入、以二维伤损热图为主输出的建模目标，并分析了纯二维监督学习在该类问题中容易出现的捷径学习与泛化不足现象。

在此基础上，本文设计了共享条件特征驱动的深度算子网络 SCF-DeepONet。该网络通过共享观测编码模块提取多激励观测下的公共表征，并分别构建二维热图分支与三维差场分支，实现了二维伤损检测主任务与三维磁矢势差场辅助任务的统一建模。随后，针对不同任务的优化特点，设计了包含热图损失、场重构损失及其阶段性组合的多任务目标函数，并提出了由热图预训练、热图--场联合训练、热图精修和 PDE 弱物理微调构成的四阶段训练策略。

总体而言，本章从问题定义、网络结构、损失函数到训练流程，系统建立了面向道岔伤损检测的 SCF-DeepONet 方法框架，为下一章开展模型优化、辅助任务作用分析及 PDE 物理约束实验奠定了基础。
